\documentclass[11pt,a4paper]{article}
\usepackage{harmonikabau_preamble}

% == Document Metadata ====================
\newcommand{\hbdoknr}{0018}
\newcommand{\hbtitel}{Audibility of Bending Modes}
\newcommand{\hbuntertitel}{Chamber Notch Filter, Transients and Torsion}
\newcommand{\hbheadtext}{Doc.\,0018 --- Audibility of Bending Modes}
\newcommand{\hbfoottext}{Cross-references: Doc.\,0016 $\cdot$ Doc.\,0017}

\AtBeginDocument{%
  \fancyhead[L]{\small\hbheadtext}%
  \fancyhead[R]{\small\thepage}%
  \fancyfoot[C]{\small\color{hb@darkgray}\hbfoottext}%
}

\begin{document}
\hbmaketitle
\noindent\textit{Correction to Doc.~0016: Bending modes can become audible under certain conditions.
Three mechanisms: chamber resonance, transient onset, torsion.
Reference: Doc.~0009, 0011, 0015--0017.}

\noindent\textit{\small Calculation script:
\textbf{pruefskript\_0018\_hoerbarkeit.py} (modal participation,
chamber notch filter, transients, torsional modes) in the same directory on GitHub.}

\section{Correction: Bending Modes Are Not Always Inaudible}

Doc.~0016 showed that the inharmonic bending modes are normally
30--50\,dB quieter than the harmonic overtones.
In practice however: with a tuneable resonance chamber, Mode~2 can be excited
in some reeds so strongly that \textbf{only Mode~2 is audible}.
And during the transient onset, Mode~2 shapes the tonal character.

\begin{tcolorbox}[colframe=red!60!black, colback=red!5!white]
\textbf{The fundamental statement stands:} Harmonic overtones dominate
in the steady state. But Mode~2 is not negligible.
\end{tcolorbox}

\section{Chamber as Notch Filter: Mode 2 Becomes Audible by Subtraction}

The chamber acts as a \textbf{notch filter} (band-rejection, Doc.~0005--0008) ---
it \textbf{suppresses} frequencies at $f_H$. The chamber does not amplify Mode~2.
Mode~2 becomes audible because the fundamental and its harmonics are attenuated,
while Mode~2 lies \textbf{outside the notch}.

Conditions: (a) $f_H$ hits a harmonic of $f_1$ $\to$ fundamental is attenuated.
(b) $f_2$ (bending mode) does \textbf{not} fall in the notch
$\to$ Mode~2 survives. This is subtraction, not addition.

\textbf{Why does this not occur in every reed?} Profiling determines $f_2$.
For a uniform reed, $f_2 \approx 6{.}27 \times f_1$ --- close to $6 \times f_1$ (a harmonic).
With stronger profiling ($f_2 \approx 4$--$5 \times f_1$),
Mode~2 falls between the harmonics and survives the notch.

\begin{tcolorbox}[colframe=keygreen, colback=green!5!white]
\textbf{In the accordion this affects notes with response problems:}
Chamber coupling is strongest there, reed plates are more heavily profiled,
and the interplay of profiling and chamber geometry happens to satisfy the
condition $f_H \neq f_2$.
\end{tcolorbox}

\section{Transient Excitation: The Onset}

When a key is pressed, the bellows pressure arrives as a step function.
Every bending mode is initially excited.
Mode~2 starts at $\approx -25$\,dB and decays in $\approx 40$\,ms.

\begin{tcolorbox}[colframe=keygreen, colback=green!5!white]
\textbf{Mode~2 is audible during the transient} --- as a brief tonal colouring
in the first 50--100\,ms. Different reeds have different $f_2$
$\to$ different tonal character.
\end{tcolorbox}

\section{Torsional Vibrations}

Torsional modes lie above the bending modes: $f_{T1} \approx 14$--$22 \times f_1$.
Torsion itself is barely audible as a pitch --- but lateral deflection
leads to channel contact in bass reeds.

At 1° torsion: bass reed (W=8\,mm): $\Delta y = 0{.}07$\,mm $>$ gap 0.04\,mm
$\to$ contact possible. \textbf{Scraping should not occur in normal operation} ---
the gaps are large enough. But in \textbf{cold instruments} scraping is normal:
Metal contracts ($\approx 2$--$3$\,µm at $\Delta T = -17$\,K),
the gap narrows, tolerance decreases. As the instrument warms up, the scraping disappears.

\begin{tcolorbox}[colframe=keygreen, colback=green!5!white]
\textbf{Paradox:} The most precisely fitted reed plates --- the most expensive ---
are most susceptible to scraping in the cold. Narrow gaps (better response)
mean less tolerance for torsion. A plate with 25\,µm gap loses
12\,\% headroom from a 3\,µm contraction --- one with 50\,µm only 6\,\%.
The cheap plate with a large gap never scrapes; the expensive one with the perfect gap
scrapes in winter.
\end{tcolorbox}

\section{Summary}

\begin{enumerate}
\item \textbf{Steady state, no chamber resonance hit:} Mode~2 at $-37$\,dB ---
  not audible. Harmonics dominate.
\item \textbf{Chamber as notch filter:} The chamber attenuates harmonics.
  If $f_2$ lies outside the notch, Mode~2 survives --- audible
  by subtraction. Particularly affects notes with response problems.
\item \textbf{Transient:} Mode~2 at $-25$\,dB, duration $\approx 40$\,ms.
  Shapes tonal character.
\item \textbf{Torsion:} $f_{T1} \approx 14$--$22 \times f_1$.
  Channel contact not expected in normal operation.
  Normal in \textbf{cold instruments} --- metal contracts, gap narrows.
  Disappears as the instrument warms up.
\item \textbf{Correction:} ``Bending modes inaudible'' applies only in steady state
  without a chamber resonance hit.
\end{enumerate}

\bigskip
\textit{Bending modes are quiet --- but not silent.
The chamber can wake them, the transient lets them briefly sound,
torsion makes them audible as scraping.}


\end{document}
